\section{Introduction}




Large language models (LLMs) have demonstrated strong reasoning capabilities that translate into substantial gains on downstream applications such as mathematics and code generation \cite{guo2025deepseek, shao2024deepseekmath}. From an architectural perspective, these capabilities have historically been tied to scale, and in particular to depth. Many reasoning tasks benefit from applying more sequential computation to an input \cite{saunshi2025reasoning}, which standard transformers realize by stacking additional layers. As a result, strong reasoners tend to be very large.
Orthogonal to increasing model size, a separate line of work improves reasoning by modifying the pretraining objective itself. Rather than relying solely on next-token prediction, these approaches adopt richer objectives such as multi-token prediction (MTP), which trains the model to predict the next $k$ tokens simultaneously. By forcing the model to \textit{think ahead}, MTP has been shown to strengthen the reasoning capabilities of LLMs \cite{liu2024deepseek, gloeckle2024better}. However, the benefits of MTP have so far been demonstrated only at larger scale.

Yet access to such models is usually constrained: frontier models are often available only via API, and many organizations 
cannot host
competitive open-weight models locally. In domains such as healthcare, this is a hard requirement; sensitive data cannot leave the premises, so models must run on-site, often on limited hardware \cite{garg2026rise}. Strong reasoning is thus increasingly needed under a parameter-efficient budget.

Looped transformers have emerged as a promising research avenue toward this goal. Rather than adding parameters, they reuse the same stack of transformer blocks across several iterations, a form of \textit{latent reasoning} in which the model refines its representation over several loops before committing to a token. This iterative reuse of computation blocks simulates the effective depth of a much larger model at a fixed parameter count, and has been shown to match the reasoning capabilities of significantly larger models \cite{fu2026simply}.

Despite their success, current looped models suffer from the following limitations. First, they overwrite intermediate computation: each iteration replaces the previous iteration's hidden representation, discarding information that may still be useful \cite{jeddi2026loopformer, zhu2025scaling}. This gives rise to \textit{latent overthinking} \cite{fu2025think}, where predictions that are already correct after an early pass are revised into errors as looping continues. Second, they suffer from \textit{undifferentiated computation}: as the number of iterations grows, hidden representations become increasingly similar across iterations, so successive iterations perform redundant work and waste compute \cite{yu2025mesh}.

To address these limitations, we (1) propose \textsc{LoopMTP}, a novel looped transformer architecture that leverages MTP in latent space as a natural remedy for latent overthinking and undifferentiated computation, (2) study the effect of the MTP alignment objective and the aggregation mechanism on downstream performance, (3) demonstrate that our approach outperforms current looped transformers, and (4) show how our architecture can be used to train small domain-specific expert models.





    
    
    
